\newtoggle{withbonus} \toggletrue{withbonus} % set true to show bonus points in grade table

% setting underlying course name (Grundlagenfach Mathematik, §-Unterricht, \ldots)
\newcommand{\mycourse}{Ergänzungsfach Informatik}
% name of the class
\newcommand{\myclass}{4. Klassen}
% set date
\newcommand{\mydate}{11.05.2026}

\title{\textbf{Komplexitätstheorie}}
\author{Thomas Graf\\
	\mycourse, \myclass}
\date{\mydate}
\maketitle
\thispagestyle{headandfoot}

\begin{center}
	\fbox{\parbox{140mm}{\centering
			\begin{itemize}
				\item Dauer der Prüfung: $45$ Minuten
				      %\item \textbf{Zeige in jeder Aufgabe unbedingt Deine Schritte!}
				\item Erlaubte Hilfsmittel:
				      \begin{itemize}
					      \item Schreibzeug
					            % \item Taschenrechner
				      \end{itemize}
			\end{itemize}}}
\end{center}
\vspace{10mm}

\begin{center}
	Vorname: \rule{45mm}{0.8pt}\qquad Nachname: \rule{45mm}{0.8pt}
\end{center}
\vspace{20mm}

\begin{center}
	\iftoggle{withbonus} {
		\combinedgradetable[h][questions]
	} {
		\gradetable[h][questions]
	}
\end{center}
\vspace{15mm}
\begin{center}
	Note:
	\vspace{10mm}

	\rule{8mm}{1.2pt}
\end{center}

\clearpage


\begin{questions}

	\question
	Begründe oder widerlege:
	\begin{parts}
        \part[1] $\lr{\frac{4}{n} + \frac{7}{n^2}} \in O\lr{\frac{1}{n}}$
        \begin{solution}
            Es gilt
            \begin{align*}
                \limsup_{n\to\infty}
                \frac{\abs{\frac{4}{n}+\frac{7}{n^2}}}{\abs{\frac{1}{n}}}
                &=
                \limsup_{n\to\infty}
                n\lr{\frac{4}{n}+\frac{7}{n^2}} \\
                &=
                \limsup_{n\to\infty}
                \lr{4+\frac{7}{n}} \\
                &=4<\infty.
            \end{align*}
            Die Behauptung ist wahr.
        \end{solution}

        \part[1] $\lr{n^3+5n^2+1} \in O\lr{n^2}$
        \begin{solution}
            Es gilt
            \begin{align*}
                \limsup_{n\to\infty}
                \frac{\abs{n^3+5n^2+1}}{\abs{n^2}}
                &=
                \limsup_{n\to\infty}
                \lr{n+5+\frac{1}{n^2}} \\
                &=\infty.
            \end{align*}
            Die Behauptung ist falsch.
        \end{solution}

        \part[1] $\lr{2^n+n^5} \in O\lr{2^n}$
        \begin{solution}
            Es gilt
            \begin{align*}
                \limsup_{n\to\infty}
                \frac{\abs{2^n+n^5}}{\abs{2^n}}
                &=
                \limsup_{n\to\infty}
                \lr{1+\frac{n^5}{2^n}} \\
                &=1<\infty.
            \end{align*}
            Da $2^n$ schneller wächst als jedes Polynom, ist die Behauptung wahr.
        \end{solution}

        \part[1] $\lr{4^n} \in O\lr{3^n}$
        \begin{solution}
            Es gilt
            \begin{align*}
                \limsup_{n\to\infty}
                \frac{\abs{4^n}}{\abs{3^n}}
                &=
                \limsup_{n\to\infty}
                \lr{\frac{4}{3}}^n \\
                &=\infty.
            \end{align*}
            Die Behauptung ist falsch.
        \end{solution}

        \part[1] $\lr{\sqrt{n}+n\log(n)} \in O\lr{n\log(n)}$
        \begin{solution}
            Es gilt
            \begin{align*}
                \limsup_{n\to\infty}
                \frac{\abs{\sqrt{n}+n\log(n)}}{\abs{n\log(n)}}
                &=
                \limsup_{n\to\infty}
                \lr{
                    \frac{1}{\sqrt{n}\log(n)}+1
                } \\
                &=1<\infty.
            \end{align*}
            Die Behauptung ist wahr.
        \end{solution}

		\droptotalpoints
	\end{parts}


    \question Gegeben sei der folgende Suchalgorithmus:
    \begin{lstlisting}[language=Python]
def count_until_greater(arr, n, limit):
    count = 0
    for i in range(n):
        count += 1
        if arr[i] > limit:
            return count
    return count
    \end{lstlisting}
    Der Algorithmus durchsucht ein Array von links nach rechts und bricht ab, sobald ein Eintrag grösser als \lstinline|limit| gefunden wird.

    \begin{parts}
        \part[1] Bestimme die Laufzeit im best-case.
        \begin{solution}
            Im best-case ist bereits der erste Eintrag grösser als \lstinline|limit|. Dann wird die Schleife nur einmal ausgeführt. Die Laufzeit ist $\Theta(1)$.
        \end{solution}

        \part[1] Die Laufzeit des Algorithmus im worst-case ist:
        \begin{checkboxes}
            \choice $\Theta(1)$
            \correctchoice $\Theta(n)$
            \choice $\Theta(n^2)$
            \choice $\Theta(\log(n))$
            \choice keine der obigen Antworten ist richtig
        \end{checkboxes}
    \end{parts} 
    \droptotalpoints


    \question Bestimme jeweils die Laufzeit der folgenden Algorithmen in Abhängigkeit von $n\in\N$.
    \begin{parts}

        \part[1] \leavevmode
        \begin{lstlisting}[language=Python]
result = 0
for i in range(n):
    for j in range(n):
        result += i + j
        \end{lstlisting}
        \begin{checkboxes}
            \choice $\Theta(n)$
            \choice $\Theta(\sqrt{n})$
            \choice $\Theta(\log(n))$
            \correctchoice $\Theta(n^2)$
            \choice $\Theta(n \log(n))$
            \choice keine der obigen Antworten ist richtig
        \end{checkboxes}
        \begin{solution}
            Beide Schleifen laufen jeweils $n$-mal. Insgesamt gibt es also $n\cdot n=n^2$ Durchläufe. Die Laufzeit ist $\Theta(n^2)$.
        \end{solution}

        \part[1] \leavevmode
        \begin{lstlisting}[language=Python]
result = 0
i = n
while i >= 1:
    result += i
    i //= 3
        \end{lstlisting}
        \begin{checkboxes}
            \choice $\Theta(1)$
            \choice $\Theta(\sqrt{n})$
            \choice $\Theta(n)$
            \choice $\Theta(n^2)$
            \correctchoice $\Theta(\log(n))$
            \choice keine der obigen Antworten ist richtig
        \end{checkboxes}
        \begin{solution}
            Die Variable $i$ wird in jedem Schritt durch $3$ geteilt. Daher gibt es ungefähr $\log_3(n)$ Durchläufe. Die Laufzeit ist $\Theta(\log(n))$.
        \end{solution}

        \part[1] \leavevmode
        \begin{lstlisting}[language=Python]
result = 0
for i in range(n):
    j = 1
    while j <= n:
        result += 1
        j *= 2
        \end{lstlisting}
        \begin{checkboxes}
            \choice $\Theta(n)$
            \choice $\Theta(\sqrt{n})$
            \choice $\Theta(\log(n))$
            \choice $\Theta(n^2)$
            \correctchoice $\Theta(n \log(n))$
            \choice keine der obigen Antworten ist richtig
        \end{checkboxes}
        \begin{solution}
            Die äussere Schleife läuft $n$-mal. Die innere Schleife verdoppelt $j$ in jedem Schritt und läuft deshalb $\Theta(\log(n))$-mal. Insgesamt ergibt sich $\Theta(n\log(n))$.
        \end{solution}

    \end{parts} 
    \droptotalpoints


    \question Der Algorithmus \lstinline|selection_sort| ist gegeben durch:
    \begin{lstlisting}[language=Python]
def selection_sort(arr):
    n = len(arr)
    for i in range(n):
        min_index = i
        for j in range(i + 1, n):
            if arr[j] < arr[min_index]:
                min_index = j
        arr[i], arr[min_index] = arr[min_index], arr[i]
    \end{lstlisting}

    \begin{parts}
        \part[2] Bestimme die worst-case Laufzeit von \lstinline|selection_sort|. Zeige deine Schritte.
        \begin{solution}
            Die äussere Schleife läuft $n$-mal. Für jedes $i$ läuft die innere Schleife von $i+1$ bis $n-1$, also ungefähr $n-i-1$-mal.

            Insgesamt ergibt sich
            \begin{align*}
                \sum_{i=0}^{n-1} (n-i-1)
                &=
                (n-1)+(n-2)+\dots+1+0 \\
                &=
                \frac{n(n-1)}{2}
                \in \Theta(n^2).
            \end{align*}

            Die worst-case Laufzeit ist also $\Theta(n^2)$.
        \end{solution}

        \part[2] Bestimme ebenso die best-case Laufzeit von \lstinline|selection_sort|. Erkläre deine Überlegungen.
        \begin{solution}
            Auch im best-case muss für jede Position $i$ das restliche Array durchsucht werden, um das Minimum zu finden. Die innere Schleife läuft also gleich oft wie im worst-case.

            Daher ist auch die best-case Laufzeit
            \[
                \Theta(n^2).
            \]
        \end{solution}
    \end{parts}
    \droptotalpoints


    \question[2] Analysiere die Laufzeit des folgenden Algorithmus:
    \begin{lstlisting}[language=Python]
def strange_sum(n):
    result = 0
    i = 1
    while i <= n:
        for j in range(n):
            result += 1
        i *= 2
    return result
    \end{lstlisting}

    Bestimme die Laufzeit des Algorithmus. Erkläre deine Überlegungen.
    \begin{solution}
        Die while-Schleife verdoppelt $i$ in jedem Schritt. Deshalb läuft sie $\Theta(\log(n))$-mal.

        Für jeden Durchlauf der while-Schleife läuft die for-Schleife $n$-mal.

        Insgesamt ergibt sich daher
        \[
            \Theta(\log(n))\cdot \Theta(n)
            =
            \Theta(n\log(n)).
        \]

        Die Laufzeit ist also $\Theta(n\log(n))$.
    \end{solution}


	\clearpage
	\begin{center}
	Es seien $f$ und $g$ Folgen reeller Zahlen.\\
    \end{center}
	\vspace{0.5cm}
	\centering
	\small
	\begin{tabular}{|c|c|c|}
		\hline
		\textbf{Notation} & \textbf{Definition}                                              & \textbf{Bedeutung}                                  \\
		\hline
		$f \in o(g)$      & $\lim\limits_{n \to \infty} \abs{\frac{f(n)}{g(n)}} = 0$         & $f$ wächst langsamer als $g$                        \\
		\hline
		$f \in O(g)$      & $\limsup\limits_{n \to \infty} \abs{\frac{f(n)}{g(n)}} < \infty$ & $f$ wächst höchstens so schnell wie $g$             \\
		\hline
		$f \in \Omega(g)$ & $\liminf\limits_{n \to \infty} \abs{\frac{f(n)}{g(n)}} > 0$      & \makecell{$f$ wächst mindestens so schnell wie $g$, \\$g\in O(f)$}                          \\
		\hline
		$f \in \Theta(g)$ & \makecell{$f \in O(g)\cap \Omega(g)$,                                                                                  \\ $0 < \liminf\limits_{n \to \infty} \abs{\frac{f(n)}{g(n)}} \leq$ \\ $\leq\limsup\limits_{n \to \infty} \abs{\frac{f(n)}{g(n)}} < \infty$} & \makecell{$f$ wächst gleich schnell wie $g$,        \\sowohl $f\in O(g)$ als auch $f\in\Omega(g)$} \\
		\hline
		$f \in \omega(g)$ & $\lim\limits_{n \to \infty} \abs{\frac{f(n)}{g(n)}} = \infty$    & \makecell{$f$ wächst schneller als $g$,             \\$g\in o(f)$}                                      \\
		\hline
	\end{tabular}

	(Wir werden meist nur $\lim$ anstelle von $\limsup$ oder $\liminf$ betrachten müssen.)


\end{questions}